\documentclass[11pt]{article}
\usepackage{booktabs, graphicx, hyperref, amsmath, geometry, array}
\geometry{margin=1in}

\title{Audio Quality Verification in AI-Generated Video: \\
A Multi-Timescale Evaluation Framework}
\author{Philo Labs}
\date{April 2026}

\begin{document}
\maketitle

\begin{abstract}
AI-generated video has reached visual photorealism, but audio quality lags far behind.
We present a corruption-based benchmark framework that injects known audio failures at
controlled severity levels and measures detection accuracy across two judge systems: a
signal processing pipeline and Gemini~2.5~Flash via OpenRouter. Across 1,094 variants
spanning four seed tasks and two real-video datasets, the judges exhibit sharply
complementary competency profiles: signal tools dominate on millisecond-precise sync
detection (99.5\% vs.\ 63.1\%), while the VLM leads substantially on artifact and SFX
perception tasks (70.6\% vs.\ 36.2\%). We present the failure taxonomy, seed task
specifications, expansion methodology, rubric design with explicit level definitions, and
experimental findings that motivate an ensemble routing strategy and a tiered reward
signal for training audio-aware video generation models.
\end{abstract}

\section{Introduction: Why Audio Verification is the Critical Bottleneck}

Visual fidelity in AI-generated video has reached the point where casual viewers often
cannot distinguish generated from real footage. Audio quality has not kept pace. Three
categories of failure are pervasive across current frontier systems.

\textbf{Voice failures.} Sora-generated dialogue routinely exhibits phoneme-timing offsets
of 200--400\,ms between speech audio and lip motion---beyond the $\sim$40\,ms perceptual
threshold---making the video look dubbed. Kling output has been observed where a
speaker's voice identity shifts mid-sentence (different timbre, different F0 range)
because the TTS module receives an ambiguous speaker embedding at a cut boundary.
Veo clips with overlapping speakers show background audio that drowns dialogue by 10--20\,dB,
with no dynamic range management.

\textbf{Sound effect failures.} Impact sounds fire 300--800\,ms after the visual contact
event (a door closing, a glass breaking), breaking the auditory-visual causality illusion.
Generated SFX frequently have incorrect spectral signatures: a wooden impact sounds
metallic, or a soft cloth rustle has the frequency profile of a hard tap. SFX are also
simply absent: a generated car collision scene may have dialogue and music but no crash
sound at all, because the model treats SFX generation as optional.

\textbf{Music failures.} The most structurally distinctive failure: music cues in AI video
\emph{do not follow shot boundaries}. A music phrase may begin 2 seconds before a cut and
continue 3 seconds past it, ignoring the scene transition entirely. This requires
sequence-level evaluation---per-clip or per-frame scoring cannot detect that the emotional
arc of a 30-second music segment is incoherent with the scene it accompanies. Additionally,
generated music often mismatches the scene's emotional valence entirely: tense orchestral
music over a comedic scene, or upbeat pop over a funeral.

No systematic infrastructure exists to detect, classify, or score any of these failures
automatically. The consequence is that generative video models receive no structured audio
feedback during training, only implicit human preference signals that conflate visual and
audio quality. This paper describes evaluation infrastructure designed to fill that gap.

\section{Seed Tasks and Ground Truth Strategy}

We define five seed tasks spanning all three audio pillars. Each seed has a mathematically
specified ground truth, enabling automated evaluation without human annotation for
Tier~1--2 tasks.

\textbf{S1 (Voice / Sync):} Inject a known time offset $\delta \in \{50, 100, \ldots,
2000\}\,\text{ms}$ between audio and video tracks. \emph{Ground truth:} offset\_ms $\neq 0$
is a failure; magnitude determines severity. Programmatically verifiable via cross-correlation.

\textbf{S2 (Voice / Speaker):} At a known swap point $t_s$, replace speech with a
frequency-shifted version (different fundamental frequency and timbre). \emph{Ground truth:}
replacement frequency ratio $\neq 1.0$ is a failure. Programmatically verifiable via
multi-feature spectral comparison (MFCC, dominant frequency ratio, spectral centroid).

\textbf{S3 (SFX / Timing):} Shift impact sound event onsets by $\delta \in \{50, 100, 200,
500, 1000, 2000\}\,\text{ms}$ relative to the visual contact event. \emph{Ground truth:}
shift\_ms $\neq 0$ is a failure. Programmatically verifiable via onset delta matching, with
VLM fallback for complex onset structures.

\textbf{S4 (SFX / Artifacts):} Inject energy spikes ($>$4$\sigma$), spectral flux anomalies
($>$4.5$\sigma$), or signal dropouts at controlled severity $s \in [0.05, 1.0]$.
\emph{Ground truth:} severity $> 0$ is a failure. Programmatically verifiable via Z-score
anomaly detection, though perceptual calibration benefits from VLM consensus.

\textbf{S5 (Music / Mood):} Replace the background music track with a track whose
valence/energy vector is displaced by mood distance $d \in \{0.2, 0.4, 0.6, 0.8, 1.0\}$
in the unit square. \emph{Ground truth:} mood\_distance $> 0.2$ is a perceptible mismatch.
Requires VLM or human judgment; programmatic proxy is the valence/energy Euclidean distance.

\subsection{Automation Boundary}

The three-tier automation model determines which evaluation layer handles each dimension:

\begin{center}
\begin{tabular}{p{3.5cm}p{3.5cm}p{4cm}}
\toprule
Signal Processing & VLM Judge & Human Required \\
\midrule
A/V sync offset (cross-corr.) &
Artifact perceptibility &
Genre-specific intent (avant-garde) \\
Onset timing delta &
Mood appropriateness &
Diegetic vs.\ non-diegetic ambiguity \\
Spectral anomaly Z-score &
Speaker naturalness &
Intentional silence vs.\ dropout \\
Frequency ratio shift &
SFX semantic match &
Cultural/contextual coherence \\
\bottomrule
\end{tabular}
\end{center}

For Tier~3 tasks, \emph{agentic pre-labeling} bridges the gap: a strong VLM annotates
semantic relationships (``the music valence is joyful but the scene shows grief''), a human
reviews a 5--10\% sample to set a calibration threshold, and the agreed labels become stable
benchmark entries. This reduces human annotation cost by approximately $10\times$ compared
to fully manual labeling.

\section{Variant Expansion and Difficulty Calibration}

From five seed tasks, we expand to 500+ variants using five systematic methods:

\textbf{1. Severity scaling.} Each continuous corruption parameter is swept across a
pre-defined grid. For S1: offset\_ms $\in \{50, 100, 150, 200, 300, 500, 750, 1000, 1500,
2000\}$ (10 levels). For S4: severity $\in \{0.05, 0.1, 0.15, 0.2, 0.3, 0.4, 0.5, 0.7,
1.0\}$ (9 levels). Combined with multiple source clips, severity scaling alone produces
100--500 variants per seed.

\textbf{2. Temporal shifting.} Varies the \emph{position} of the corruption within the clip,
not just its magnitude. A 200\,ms sync offset applied at second 1.0 is harder to detect than
the same offset at second 6.0 because early-clip onset ambiguity provides less comparison
context for cross-correlation.

\textbf{3. SNR manipulation.} Controls the signal-to-noise ratio between the injected
corruption and the background audio. An artifact buried under a music track at SNR $= -3$\,dB
is nearly imperceptible to both signal tools and VLMs. SNR levels $\in \{20, 10, 5, 0, -3\}$\,dB
expand S4 into a masking-difficulty axis.

\textbf{4. Distractor injection.} Adds a competing audio event temporally close to the
target event. For S3 (SFX timing), a distractor impact sound 100\,ms before the target
confuses both onset detectors and human listeners, creating hard negatives.

\textbf{5. Cross-pillar combination.} Applies two corruptions simultaneously (e.g., sync
drift + artifact injection). Composite variants test whether judges can disentangle
co-occurring failures and are the most realistic failure mode in actual generated video.

\subsection{Difficulty Calibration}

A useful benchmark requires tasks where current models fall in the 35--75\% pass-rate band:
wide enough to distinguish strong from weak models, narrow enough that improvement is
measurable. Our framework categorizes variants as \emph{easy} ($>$75\%), \emph{medium}
(35--75\%), or \emph{hard} ($<$35\%) after each run and flags any seed where the overall
pass rate is outside the band. In our experiments, S1 is at ceiling for signal (99.5\%)
and needs adversarial distractor variants to become discriminating; S2 and S3 are below
the floor for both judges and need easier severity levels (larger frequency shifts, longer
temporal offsets) to bring them into the band.

\section{Verifiable Rubric Design}

Five scoring dimensions, each on a 1--5 scale with explicit level definitions:

\textbf{A/V Synchronization} measures audio-to-visual temporal alignment.
1 = catastrophic offset ($>$500\,ms), clearly dubbed or disjointed;
2 = severe offset (200--500\,ms), distracting to any viewer;
3 = noticeable offset (80--200\,ms), perceptible on focused viewing;
4 = minor offset (40--80\,ms), at perceptual threshold;
5 = imperceptible offset ($<$40\,ms). \emph{Automation:} fully signal-verifiable.

\textbf{Artifact Quality} measures freedom from introduced noise, clipping, or dropouts.
1 = severe distortion or long dropouts dominate the clip;
2 = prominent clicks, pops, or clipping that interrupt listening;
3 = audible artifacts but tolerable;
4 = subtle artifacts detectable on headphones;
5 = artifact-free. \emph{Automation:} Z-score detection for severe cases; VLM for subtle.

\textbf{Speaker Consistency} measures voice identity stability across clip duration.
1 = completely different voice (different gender, age, or accent);
2 = notable timbre or pitch change inconsistent with natural variation;
3 = slight inconsistency that attentive listeners notice;
4 = minor variation within plausible natural range;
5 = fully consistent speaker identity. \emph{Automation:} spectral composite for extremes; VLM for subtle.

\textbf{Semantic Match} measures appropriateness of SFX selection for the visual event.
1 = completely wrong (metallic sound for cloth, animal sound for mechanical action);
2 = major mismatch (wrong category);
3 = plausible but imprecise (wood when plastic expected);
4 = mostly appropriate with minor spectral discrepancy;
5 = perfectly matched to visual cause. \emph{Automation:} VLM required for all levels.

\textbf{Music Coherence} measures alignment between music's emotional valence/energy and scene mood.
1 = directly opposite valence (triumphant music over tragedy);
2 = clearly inappropriate for scene type;
3 = neutral/generic music with no clear mismatch but no enhancement;
4 = generally fitting with minor mood incongruity;
5 = enhances narrative with well-matched dynamics and valence. \emph{Automation:} VLM required; human review for borderline cases.

\subsection{The Cross-Cut Music Challenge}

The spec identifies a structural property that makes music evaluation uniquely hard: music
cues span shot boundaries. A musical phrase beginning two seconds before a cut and ending
three seconds after it cannot be evaluated at the clip level. Per-clip scoring would assign
the music a neutral or positive score (``the music matches this 8-second window'') while
missing that the phrase as a whole ignores the emotional transition at the cut point.

Correct sequence-level evaluation requires maintaining a \emph{music state buffer} across
consecutive clips: tracking phrase position, key, tempo, and valence over 20--60 second
windows that span multiple cuts. The rubric's Music Coherence dimension is designed to be
applied at this multi-clip granularity, not per-clip. Our current benchmark evaluates
8-second clips, which cannot capture cross-cut coherence. Sequence-level evaluation is
identified as the primary architectural extension required for full music pillar coverage.

\subsection{Rubric Edge Cases}

Three scenarios where the rubric struggles reveal fundamental ambiguities in the ground
truth definition:

\textbf{1. Intentional silence.} A 3-second dramatic pause in dialogue is not a signal
dropout, but our S4 artifact detector scores it as severity~1.0 because energy drops to
$-$60\,dBFS. The rubric's Artifact Quality dimension requires context: is the silence
preceded by a charged emotional moment (intentional) or does it occur mid-sentence
(malfunction)? No signal-only system can make this determination.

\textbf{2. Diegetic vs.\ non-diegetic audio.} A scene showing a character listening to a
radio contains both diegetic audio (the radio) and non-diegetic audio (the film score).
Mood Coherence dimension scoring is ambiguous: the diegetic radio may play music that
mismatches the scene's emotional tone as a \emph{narrative choice} (ironic counterpoint),
not a generation failure. The rubric needs a diegetic flag to handle this correctly.

\textbf{3. Genre-stylistic mismatch.} Horror and experimental film deliberately use
out-of-sync audio, uncanny sound effects, and dissonant music. Our benchmark would flag
a correctly generated horror clip as failing on A/V Synchronization and Music Coherence.
Genre metadata must gate the rubric: the acceptable range for each dimension shifts
substantially across genre categories.

\section{Signal Processing vs.\ Model-Based Detection}

\subsection{Judge Architectures}

The \textbf{signal processing pipeline} uses librosa for FFT cross-correlation (S1), a
weighted multi-feature spectral composite (S2), onset delta matching (S3), and Z-score
anomaly detection (S4). Its defining characteristic is \emph{deterministic precision}: given
identical inputs, it produces identical outputs, and its error mode is a hard miss rather
than a soft hallucination.

The \textbf{VLM judge} (Gemini~2.5 Flash via OpenRouter) receives the full remuxed video
(corrupted audio muxed back onto the source video track) as a base64-encoded MP4, enabling
joint audio-visual reasoning. For each variant, it returns a structured JSON response with
per-dimension rubric scores plus a binary detection flag. Its defining characteristic is
\emph{perceptual holism}: it reasons about audio quality as a viewer, not as a signal
analyst, which makes it robust to stimuli that lack clean spectral signatures but brittle
to precise temporal measurement.

\subsection{Experimental Results}

We evaluated four seed tasks across 1,094 variants on AVA \cite{gu2018ava} and Greatest
Hits \cite{owens2016visually} datasets using real API calls to Gemini~2.5~Flash via
OpenRouter alongside the signal processing pipeline.

\begin{center}
\begin{tabular}{lrrrr}
\toprule
Seed & Variants & Signal Acc. & VLM Acc. & Leading Judge \\
\midrule
S1 (sync, AVA)          & 198 & 99.5\% & 63.1\% & Signal ($+$36.4\,pp) \\
S2 (speaker, AVA)       & 198 & 33.3\% & 43.9\% & VLM ($+$10.6\,pp) \\
S3 (SFX timing, GH)     & 225 & 11.1\% & 31.1\% & VLM ($+$20.0\,pp) \\
S4 (artifacts, AVA)     & 473 & 36.2\% & 70.6\% & VLM ($+$34.4\,pp) \\
\bottomrule
\end{tabular}
\end{center}

\textbf{S1} validates the signal pipeline's core competency: FFT cross-correlation achieves
near-perfect sync detection across all difficulty levels. The VLM achieves only 54\% on hard
variants ($<$300\,ms offset), confirming that language-model temporal reasoning carries
$\sim$200\,ms noise that is unavoidable without explicit signal computation.

\textbf{S2} exposes a fundamental weakness in spectral feature-matching on real speech: MFCC
cosine similarity between original and frequency-shifted replacement remains above 0.95 on
AVA clips, making all spectral features nearly useless. The VLM's marginal advantage (+10.6\,pp)
suggests it captures suprasegmental cues (rhythm, intonation envelope) that our hand-crafted
feature composite misses---pointing toward learned speaker embeddings as the correct
signal-side approach.

\textbf{S3} produces the most surprising result: signal achieves only 11.1\% on the Greatest
Hits dataset, far below the VLM's 31.1\%. Both are below the discriminating band. The failure
source is compound onset structure in real material-impact sounds---a wooden strike has an
initial contact transient followed by a resonance buildup, and our onset delta-matching logic
confuses the two phases. This demonstrates that ``Tier~1 by construction'' (known ground truth)
does not imply ``signal-dominated in practice.'' The appropriate detection method depends on
acoustic complexity, not just ground truth structure.

\textbf{S4} shows the VLM dominating by 34.4\,pp. Artifact detection is a perceptual judgment:
trained listeners recognize distortion and spectral tearing holistically, without needing a
reference signal. The pipeline's Z-score approach requires a clean baseline distribution estimated
from the signal itself, which is poorly calibrated on 8-second clips where artifacts occupy
a significant fraction of the clip. The VLM's zero-shot perceptual query (``does this sound
degraded?'') requires no baseline. This is the domain where holistic reasoning is the correct
computational model. Across 1,094 variants, 137 cases favor signal only and 294 favor VLM only,
confirming genuine complementarity.

\section{Toward Automated Task Generation}

Our corruption-based framework provides a scalable path to benchmark growth with minimal human
input. The key design choices that enable this are: (1) separating \emph{corruption application}
from \emph{ground truth extraction}, so new datasets can be onboarded by providing only source
clips; (2) SHA-256-based deterministic seeding, ensuring that the same source clip and parameter
set always produce the same variant and preventing evaluation pollution from randomness; (3)
automatic difficulty calibration after each run, which drives the expansion strategy by flagging
which severity ranges need more coverage; and (4) the agentic pre-labeling workflow for Tier~3
tasks, which converts expensive human judgment into a one-time calibration cost.

From five seed tasks, the framework generates 500+ variants without new task design effort. The
next scaling step---compositional corruption (sync drift + artifact, speaker swap + SFX mistiming)
---requires no new seed definitions, only a parameter combining step in the expansion module.
The full automation limit is Tier~3 mood coherence at the sequence level, which requires the
music state buffer architecture described in Section~5.2.

\section{Rubric as Reward Signal for Training}

The five rubric dimensions map directly to a RLVR reward structure. Tier~1 dimensions (A/V
Synchronization, Artifact Quality) provide \emph{hard verifiable rewards}: a signal tool
measures the offset and assigns a score deterministically. A sync offset of 200\,ms yields
score~2; improving to 80\,ms yields score~4. This is the same structure as code execution
feedback in reasoning model training: the reward is not a preference estimate but a
measurement.

Tier~3 dimensions (Semantic Match, Music Coherence) provide \emph{soft verifiable rewards}
via VLM consensus. Three independent VLM queries scoring the same clip, with the majority
vote as the reward signal, achieve sufficient consistency for training on cases with $d > 0.6$
mood distance. Borderline cases ($d < 0.3$) should be excluded from the reward signal or
down-weighted, as VLM disagreement in that range is high.

The critical advantage of the tiered structure over a single holistic score is
\emph{decomposed optimization pressure}. A model with perfect sync but poor mood matching
receives a gradient that specifically penalizes Music Coherence while preserving A/V
Synchronization. A holistic score conflates these dimensions and produces ambiguous
gradients. The rubric's five-dimension decomposition thus functions simultaneously as a
structured evaluation report and as a precision loss function.

\section{Limitations}

\textbf{Gemini's sync accuracy with full visual context (63.1\%) reframes S1.}
The pipeline remuxes corrupted audio back into the source video before scoring, so
Gemini received the complete video+audio stream for all AVA and Greatest Hits variants.
The 63.1\% accuracy on S1 therefore represents Gemini's genuine multimodal sync-detection
performance with lip-motion reference available. The fact that it still reaches only
54\% on hard variants ($<$300\,ms offset) demonstrates that language-model temporal
reasoning cannot reliably detect sub-perceptual-threshold offsets even with visual
reference frames---not a data-collection gap, but a fundamental capability limit of
autoregressive VLMs relative to cross-correlation. Future work should evaluate whether
chain-of-thought prompting asking the model to reason frame-by-frame improves this.

\textbf{Single-pass VLM scoring.} Each variant was scored by a single Gemini API call.
VLMs are stochastic: the same prompt and video can yield different binary detection
outputs on repeated queries, with variance estimated at $\pm$5--10\,pp at the aggregate
level. Consensus scoring---three independent queries per variant, majority vote as the
label---would reduce this noise substantially at $3\times$ API cost. Results for S2 and
S3, where aggregate accuracy is near chance, should be interpreted with this variance in mind.

\textbf{S5 (Music / Mood) not evaluated.} The fifth seed task---music mood swap with
controlled valence/energy distance---was implemented but not included in the reported
runs, because AVA and Greatest Hits clips do not carry background music tracks. Evaluating
S5 requires music-bearing clips (e.g., Condensed Movies) or synthetic mixed clips, and
constitutes the primary gap in pillar coverage: the Music pillar has no measured results.

\textbf{Sequence-level music evaluation absent.} All experiments score 8-second clips;
the Music Coherence dimension as currently implemented cannot detect cross-cut phrase
incoherence as described in Section~4.2.

\textbf{Edge cases require schema extensions.} Intentional silence, diegetic audio, and
genre-stylistic choices require genre metadata and diegetic flags that are not yet in
the benchmark schema.

\begin{thebibliography}{9}
\bibitem{gu2018ava}
R.~Gu et al., ``AVA: A Video Dataset of Spatio-temporally Localized Atomic Visual Actions,''
\emph{CVPR}, 2018.

\bibitem{owens2016visually}
A.~Owens et al., ``Visually Indicated Sounds,'' \emph{CVPR}, 2016.
\end{thebibliography}

\end{document}
